\documentclass{article}
\usepackage{hyperref}
\usepackage[margin=0.5in]{geometry}
\usepackage{amsmath}
\usepackage{graphicx}
\usepackage[bottom]{footmisc}

\graphicspath{ {./assets/} }

\title{PHYSIOL 578: HW 2.2\footnotetext{U-M GPT was used to help draft and refine the answers.}}
\author{Jack Yang \\ \href{mailto:yhyunjoo@umich.edu}{yhyunjoo@umich.edu}}

\begin{document}
\maketitle

\section{
    BDNF is an important neurotrophic factor that promotes neurogenesis and synaptogenesis.
    BDNF binds to its receptor TrkB on the plasma membrane.
    Please describe how BDNF activates TrkB and its downstream signaling pathways, and how you would experimentally test activation of these pathways.
}
There are two main signaling pathways when BDNF activates TrkB.

The first is the Ras-Raf-MEK-ERK pathway.
BDNF binds TrkB, causing receptor dimerization and tyrosine autophosphorylation.
Phosphorylated TrkB recruits SHC, followed by Grb2 and SOS.
SOS activates the small GTPase Ras.
Ras initiates the kinase cascade: Ras$\rightarrow$Raf$\rightarrow$MEK1/2$\rightarrow$ERK1/2.
ERK phosphorylates cytoplasmic and nuclear proteins, changing gene expression.

The second is the PLC$\gamma$-IP$_3$ pathway.
The activated TrkB recruits and phosphorylates PLC$\gamma$.
PLC$\gamma$ cleaves the membrane lipid PIP$_2$ into PIP$_3$ and DAG.
IP$_3$ binds IP$_3$ receptors on the ER, releasing Ca$^{2+}$ into the cytosol.
This activates calmodulin-dependent signaling, including CaMKs.
At the same time, DAG, together with Ca$^{2+}$ ctivates protein kinase C.

You could test the activation of these pathways by treating cells expressing TrkB with BDNF and comparing them with untreated controls.
Then you can use Western blots with phospho-specific antibodies to measure activation of TrkB and its downstream pathways, including p-TrkB, p-ERK1/2 for the MAPK pathway, p-AKT for the PI$_3$K-AKT pathway, and p-PLC$\gamma$ for the PLC$\gamma$ pathway, while measuring the corresponding total proteins as controls.
Co-immunoprecipitation could additionally test whether activated TrkB associates with adaptor proteins.

\end{document}